\documentclass[11pt]{article}
\usepackage[utf8]{inputenc}
\usepackage{amsmath, amssymb, amsthm}
\usepackage{geometry}
\usepackage{enumitem}
\usepackage{longtable, booktabs}
\usepackage{array}
\usepackage{hyperref}
\usepackage[dvipsnames]{xcolor}
\usepackage[most]{tcolorbox}
\hypersetup{colorlinks=true, linkcolor=NavyBlue, citecolor=ForestGreen, urlcolor=NavyBlue}
\geometry{margin=1in}

\theoremstyle{plain}
\newtheorem{theorem}{Theorem}[section]
\newtheorem{proposition}[theorem]{Proposition}
\newtheorem{corollary}[theorem]{Corollary}
\newtheorem{lemma}[theorem]{Lemma}
\theoremstyle{definition}
\newtheorem{definition}[theorem]{Definition}
\theoremstyle{remark}
\newtheorem{remark}[theorem]{Remark}

\newcommand{\Vsub}{V_{600}}
\newcommand{\pphi}{\varphi}
\newcommand{\Cphi}{C_{\pphi}}
\newcommand{\Q}{\mathbb{Q}}
\newcommand{\R}{\mathbb{R}}
\newcommand{\Z}{\mathbb{Z}}
\newcommand{\op}{\mathrm{op}}
\newcommand{\Schur}{\operatorname{Schur}}

\title{The Cascade Curvature Operator under Refinement: \\
Schur--Complement Identity and Energy Intertwining \\[0.4em]
\large A predecessor to the cascade-closure mechanism paper}
\author{Lee Smart \\ Institute of Vibrational Field Dynamics}
\date{2026-05}

\begin{document}

\maketitle

\begin{abstract}
We isolate a clean abstract model --- a scalar vertex-sector graph
Laplacian under pure edge-midpoint subdivision with degree-$2$
midpoints --- and prove a Schur-complement halving identity in this
model that lifts the source's coboundary factor-$\tfrac12$ relation
$p^{1} d_{n+1} = \tfrac12 d_n p^{0}$ from
\cite{CascadeClosureDynamics} to the curvature level. With the
scaled vertex Laplacian $\widetilde A_n := 2^n A_n^{\mathrm{vertex}}$,
the Schur complement under the (boundary, interior) partition equals
the previous-level operator exactly:
$\Schur(\widetilde A_{n+1}) = \widetilde A_n$. The identity is
verified analytically and numerically on two concrete instances of
the abstract model (single edge with one-midpoint subdivision;
triangle with edge-midpoint subdivision) and yields the
cascade-mechanism residual-monotonicity hypothesis
$\widetilde R_n(p^{0} \psi) \le \widetilde R_{n+1}(\psi)$ for every
$\psi \in X_{n+1}^{0}$, with equality if (but not in general only if)
$\psi$ is the harmonic extension of $p^{0} \psi$. Strict
operator-level refinement-compatibility
$p^{0} \widetilde A_{n+1} = \widetilde A_n p^{0}$ remains false; the
defect operator vanishes on harmonic extensions and is generically
non-zero on the level-$n$-supported sector
$\{\psi : \psi_I = 0\}$. We \emph{do not} extend the result to the
full Schl\"{a}fli-refinement P3/P4 tower of
\cite{CascadeRefinementSpaces, CascadeClosureDynamics} (which adds
face-centroids, midpoint-centroid edges, and $F$-valued chain
spaces with antisymmetric oriented-edge convention); the gap between
this paper's abstract scalar pure-midpoint model and that fuller
setting is named explicitly in \S\ref{sec:scope-and-gap} as the
remaining upstream task. The result discharges
cascade-mechanism's obligation~(O3) for the abstract model and
partially discharges (O2) at the boundary-energy / Schur level for
the abstract model, recasting the cascade-mechanism worked-invocation
status table accordingly.
\end{abstract}

\begin{tcolorbox}[title=\textbf{Scope}]
\begin{itemize}[leftmargin=*]
\item This is a narrow technical paper on an \emph{abstract scalar
  model} of refinement (vertex-sector graph Laplacian, pure
  edge-midpoint subdivision with degree-$2$ midpoints), not the full
  Schl\"{a}fli-refinement P3/P4 tower of
  \cite{CascadeRefinementSpaces, CascadeClosureDynamics} (see
  \S\ref{sec:scope-and-gap}). It contains no new mechanism, no
  programme framing, no consciousness or cosmology claim. It is the
  source-side companion to \cite{CascadeMechanism}, applied to the
  abstract model.
\item Within the abstract scalar model: (O0) closed; (O3) closed via
  the Schur halving identity (Theorem~\ref{thm:schur},
  Theorem~\ref{thm:O3-discharge}); (O2) closed at the Schur /
  boundary-energy level only (Proposition~\ref{prop:strict-fails});
  (O1) recorded but not nontrivially discharged
  (Proposition~\ref{prop:O1-status}).
\item Lifting these results to the full P3/P4 tower (centroids,
  midpoint-centroid extra edges, $F$-valued $X_n^{0}$ and
  antisymmetric $X_n^{1}$ chain spaces) is named as the remaining
  upstream task and is not attempted here.
\end{itemize}
\end{tcolorbox}

\section{Introduction and source-state recap}\label{sec:intro}

The cascade-closure mechanism paper~\cite{CascadeMechanism} defines
cascade closure as a rung-indexed closure-residual / projection-stack
event, proves a conditional residual-zero / fixed-point propagation
proposition, and supplies a consumer API listing four obligations
that successor papers must discharge to invoke cascade closure as a
load-bearing premise:
\begin{enumerate}[leftmargin=2.4em, label=\textup{(O\arabic*)}, start=0, nosep]
\item Standing analytic / admissibility hypotheses.
\item Top-rung fixed point and convergence.
\item Flow intertwining
  $\pi^{\sharp}_{k+1, k} \circ \Psi_{k+1}^{t} = \Psi_k^{t} \circ \pi^{\sharp}_{k+1, k}$.
\item Residual monotonicity
  $R_k(\pi^{\sharp}_{k+1, k} \psi) \le R_{k+1}(\psi)$.
\end{enumerate}

\cite{CascadeMechanism}~\S3.5 records the per-obligation status in
the imported P3/P4 finite-level tower of
\cite{CascadeRefinementSpaces, CascadeClosureDynamics}: (O0) is
satisfied at the finite-dimensional level; (O1)--(O3) are open in
the source state, with (O2) closed only for the source's mass-only
operator $L_n = -\gamma C_n$ at $\alpha = \beta = 0$, not for the
curvature/residual operator $A_n$ that the cascade-mechanism flow
naturally uses. \cite{CascadeMechanism}~\S3.5 names the minimal
upstream upgrade required and observes that the role of
\cite{CascadeMechanism} is to make the obligations visible as the
programme's actual dependency gate, not to discharge them.

The present paper is a partial source-side discharge, carried out
in a deliberately scaled-down setting: an \emph{abstract scalar
pure-midpoint refinement model} (Definition~\ref{def:abstract-model}
below) that retains the load-bearing coboundary refinement relation
of the imported tower while dropping the centroids,
midpoint-centroid edges, and $F$-valued / antisymmetric oriented-edge
structure of the full Schl\"{a}fli refinement. In this abstract
model, with the curvature operator $A_n$ rather than the source's
mass-only $L_n$, we supply a single source-side identity and
exhibit its consequences for (O2) and (O3). The lift to the full
P3/P4 tower of \cite{CascadeRefinementSpaces, CascadeClosureDynamics}
is named explicitly in \S\ref{sec:scope-and-gap} as hypotheses
\textup{(L1)}--\textup{(L3)} and is not attempted here. We do not
redefine the cascade-closure mechanism, do not modify
\cite{CascadeMechanism}'s consumer API, and do not introduce any
programme-level claim.

\subsection{Scope: abstract scalar pure-midpoint model vs.\ full P3/P4 tower}\label{sec:scope-and-gap}

The actual source refinement of \cite{CascadeRefinementSpaces}
(\texttt{cascade-refinement-spaces.tex} lines~273--287) is the
\emph{Schl\"{a}fli refinement}: at each level $n$ it adds, alongside
the new edge-midpoints $M_n^{\bullet}$, a separate set of
face-centroids $C_n^{\bullet}$ (one per $2$-face of the underlying
$4$-polytope), together with extra level-$(n+1)$ edges connecting
midpoints to centroids and centroids to other midpoints in a
$W^{\bullet}$-equivariant pattern. The chain spaces in the source
are also $F$-valued (\texttt{cascade-refinement-spaces.tex}
lines~464--493): $X_n^{0, \bullet}$ is $F^{0}$-valued on
$V(G_n^{\bullet})$ and $X_n^{1, \bullet}$ is $F^{1}$-valued and
antisymmetric on \emph{oriented} edges, with edge-parent map
$\pi_{n+1, n}^{E}$ taking the value $\bot$ on the new
midpoint-centroid and centroid-midpoint edges.

The present paper does \emph{not} work in this fuller setting. We
work in a deliberate \emph{abstract scalar pure-midpoint model}
defined as follows:

\begin{definition}[Abstract scalar pure-midpoint refinement model]\label{def:abstract-model}
A pair of finite undirected graphs $G_n = (V_n, E_n)$ and
$G_{n+1} = (V_{n+1}, E_{n+1})$ together with bonding maps is in the
abstract scalar pure-midpoint model when:
\begin{enumerate}[label=\textup{(M\arabic*)}, leftmargin=2.4em, nosep, topsep=2pt]
\item $V_{n+1} = V_n \sqcup M_n$, where $M_n$ contains exactly one
  new \emph{midpoint} vertex $m_e$ per edge $e \in E_n$, and no other
  new vertices (no centroids);
\item $E_{n+1} = \{(u, m_e), (m_e, v) : e = (u, v) \in E_n\}$, i.e.\
  every level-$(n+1)$ edge is a child of exactly one level-$n$ edge
  (no extra midpoint--midpoint edges);
\item the chain spaces are scalar:
  $X_n^{0} = \R^{V_n}$, $X_n^{1} = \R^{E_n}$, with standard $\ell^{2}$
  inner products and \emph{unoriented} edge convention;
\item bonding maps are vertex-restriction
  $p^{0}_{n+1, n} : X_{n+1}^{0} \to X_n^{0}$ along the inclusion
  $V_n \subset V_{n+1}$ and parent-fibre average
  $p^{1}_{n+1, n} : X_{n+1}^{1} \to X_n^{1}$ by factor $\tfrac12$ on
  each of the two children of every parent edge.
\end{enumerate}
The vertex-sector curvature is the standard graph Laplacian
$A_n^{\mathrm{vertex}} := d_n^{*} d_n$ on $X_n^{0}$, with
$d_n : X_n^{0} \to X_n^{1}$ the oriented difference (with respect to
an arbitrary but fixed orientation of $E_n$).
\end{definition}

In this abstract model, the coboundary refinement relation
\begin{equation}\label{eq:cobref}
   p^{1}_{n+1, n} d_{n+1} = \tfrac12 \, d_n \, p^{0}_{n+1, n}
\end{equation}
holds verbatim, by direct computation: for any
$f \in X_{n+1}^{0}$, $(p^{1} d_{n+1} f)(e) = \tfrac12 [(d_{n+1} f)(u, m_e)
+ (d_{n+1} f)(m_e, v)] = \tfrac12 [(f(m_e) - f(u)) + (f(v) - f(m_e))]
= \tfrac12 (f(v) - f(u)) = \tfrac12 (d_n p^{0} f)(e)$. We use
\eqref{eq:cobref} as the load-bearing input below. The relation also
holds in the source's full Schl\"{a}fli setting~(P4
\texttt{coboundary-refinement}, lines~896--937), but the present
paper only uses the abstract-model version, where it is elementary.

\paragraph*{The gap to the full P3/P4 tower.}
The Theorems and Proposition below all hold only \emph{within
Definition~\ref{def:abstract-model}'s abstract model}. They do
\emph{not} immediately transfer to the full Schl\"{a}fli setting of
\cite{CascadeRefinementSpaces}. The minimal lift hypotheses we
identify are:
\begin{enumerate}[label=\textup{(L\arabic*)}, leftmargin=2.4em, nosep, topsep=2pt]
\item the level-$(n+1)$ vertex Laplacian, restricted to boundary
  vertices $V_n \subset V(G_{n+1}^{\bullet})$, agrees with the
  Laplacian of the abstract model up to a defect that is supported
  only on midpoints (no extra contribution from centroid neighbours
  or centroid-midpoint edges);
\item the parent-fibre averaging map $p^{1}_{n+1, n}$ on
  $F^{1}$-valued, antisymmetric oriented-edge chains satisfies the
  calibration identity $p^{1}(p^{1})^{*} = \tfrac12 \, I$ in the
  $F^{1}$ inner product, with the same factor of $\tfrac12$ as in
  the scalar abstract model;
\item the harmonic-extension operator $H_{n+1}$ in the
  $F^{0}$-valued setting saturates the
  energy lower bound coming from
  \textup{(L1)} and \textup{(L2)}.
\end{enumerate}
Each of these is concrete and isolatable: (L1) is a graph-theoretic
incidence statement about the $W^{\bullet}$-equivariant midpoint
and centroid placement; (L2) is an inner-product calculation
in the source's $F^{1}$-valued, antisymmetric oriented-edge
convention~(\texttt{cascade-refinement-spaces.tex} lines~464--493);
(L3) is the analogue of Step~3 of the proof of
Theorem~\ref{thm:schur} in the $F$-valued setting. Lifting the
Schur-halving identity to the full P3/P4 tower under these
hypotheses is named here as the remaining upstream task and is not
attempted in this paper.

\subsection{Source objects (recap)}

Within the abstract model of Definition~\ref{def:abstract-model}, we
use scalar finite-dimensional chain spaces. The vertex-block
curvature operator is
\begin{equation}\label{eq:Anvertex}
   A_n^{\mathrm{vertex}} := d_n^{*} d_n \in \R^{V_n \times V_n},
\end{equation}
the standard graph Laplacian on $V_n$.

The bonding maps are the pair
$p_{n+1, n} = (p^{0}_{n+1, n}, p^{1}_{n+1, n})$, where
$p^{0}_{n+1, n} : X_{n+1}^{0} \to X_n^{0}$ is vertex-restriction
along the canonical inclusion $V_n \subset V_{n+1}$ and
$p^{1}_{n+1, n} : X_{n+1}^{1} \to X_n^{1}$ is the parent-fibre
average by factor $\tfrac12$ on each of the two children of every
parent edge.

The coboundary refinement relation \eqref{eq:cobref} of
\S\ref{sec:scope-and-gap} is the load-bearing input below. (The
analogous identity is also a theorem in the source's full
Schl\"{a}fli setting, \cite{CascadeClosureDynamics}~Proposition
\texttt{coboundary-refinement} at lines~896--937.)

\subsection{What this paper proves}

The paper's central object is the scaled vertex curvature operator:
\begin{definition}[Scaled curvature operator]\label{def:Atilde}
$\widetilde A_n := 2^n A_n^{\mathrm{vertex}}$, defined on $X_n^{0}$.
The corresponding scaled residual is
$\widetilde R_n(\phi) := \tfrac12 \langle \phi, \widetilde A_n \phi\rangle = 2^{n-1} \langle \phi, A_n^{\mathrm{vertex}} \phi\rangle$.
\end{definition}

The scaling factor $2^n$ is the unique factor that compensates the
refinement-induced halving in \eqref{eq:cobref} once one passes from
the coboundary level to the curvature ($d^{*} d$) level. The main
results are an exact Schur-complement identity for $\widetilde A_n$
under boundary--interior decomposition (Theorem~\ref{thm:schur}), a
clean energy-intertwining inequality on the harmonic-extension
sector that discharges cascade-mechanism's (O3)
(Theorem~\ref{thm:O3-discharge}), and an explicit defect identity
showing that strict operator-level refinement-compatibility for
$\widetilde A_n$ remains false (Proposition~\ref{prop:strict-fails}).

\section{Schur-complement identity}\label{sec:schur}

Fix a level $n$ and a finite refinement
$V(G_n^{\bullet}) \subset V(G_{n+1}^{\bullet})$ such that
$|V(G_{n+1}^{\bullet})| - |V(G_n^{\bullet})|$ new ``midpoint''
vertices are added (each created as the midpoint of one existing
edge of $G_n^{\bullet}$ that was subdivided to produce
$G_{n+1}^{\bullet}$). Partition the vertex set
$V(G_{n+1}^{\bullet}) = B \sqcup I$ into
\begin{itemize}[nosep]
\item $B := V(G_n^{\bullet})$, the \emph{boundary} vertices (those
  retained from level~$n$);
\item $I := V(G_{n+1}^{\bullet}) \setminus V(G_n^{\bullet})$, the
  \emph{interior} or \emph{new} vertices.
\end{itemize}
Block-decompose $A_{n+1}^{\mathrm{vertex}}$ in this $(B, I)$ ordering:
\begin{equation}\label{eq:block}
   A_{n+1}^{\mathrm{vertex}} =
     \begin{pmatrix} A_{BB} & A_{BI} \\ A_{IB} & A_{II} \end{pmatrix},
\end{equation}
where
$A_{BB} \in \R^{B \times B}$,
$A_{BI} = A_{IB}^{T} \in \R^{B \times I}$,
$A_{II} \in \R^{I \times I}$.

By construction $A_n^{\mathrm{vertex}}$ is the graph Laplacian on
$V(G_n^{\bullet})$. The harmonic extension
$H_{n+1} : X_n^{0} \to X_{n+1}^{0}$ is the unique linear map sending
$f \in X_n^{0}$ to the level-$(n+1)$ state with prescribed boundary
values $f$ and interior values
$f_I = -A_{II}^{-1} A_{IB} f$, equivalently the unique minimiser of
the level-$(n+1)$ energy
$\tfrac12 \langle \,\cdot\, , A_{n+1}^{\mathrm{vertex}}\,\cdot\,\rangle$
subject to the boundary constraint
$p^{0}_{n+1, n} \cdot = f$. We assume $A_{II}$ is invertible (the
standard finite-dimensional case in
\cite{CascadeRefinementSpaces, CascadeClosureDynamics}).

The Schur complement of $A_{n+1}^{\mathrm{vertex}}$ with respect to
the partition $(B, I)$ is
\begin{equation}\label{eq:SchurDef}
   \Schur(A_{n+1}^{\mathrm{vertex}}) :=
     A_{BB} - A_{BI} A_{II}^{-1} A_{IB}
     \in \R^{B \times B}.
\end{equation}

\begin{theorem}[Schur--complement halving in the abstract model]\label{thm:schur}
In the abstract scalar pure-midpoint refinement model of
Definition~\ref{def:abstract-model}, with vertex-sector curvature
$A_n^{\mathrm{vertex}} = d_n^{*} d_n$, the Schur complement of
$A_{n+1}^{\mathrm{vertex}}$ under the (boundary, interior) partition
\eqref{eq:block} is exactly half the level-$n$ Laplacian:
\begin{equation}\label{eq:schur-half}
   \Schur(A_{n+1}^{\mathrm{vertex}}) =
     \tfrac12 \, A_n^{\mathrm{vertex}}.
\end{equation}
Equivalently, with $\widetilde A_n := 2^n A_n^{\mathrm{vertex}}$:
\begin{equation}\label{eq:schur-tilde}
   \Schur(\widetilde A_{n+1}) = \widetilde A_n.
\end{equation}
\end{theorem}

\begin{proof}
We proceed in three steps: identify the Schur complement with a
harmonic energy, decompose that energy via the source's coboundary
refinement relation, and verify directly that the harmonic
extension saturates the resulting lower bound.

\smallskip
\textbf{Step 1.}
For any positive-semi-definite $A_{n+1}^{\mathrm{vertex}}$ with
invertible interior block $A_{II}$, the Schur complement and the
harmonic-extension energy agree:
$\langle f, \Schur(A_{n+1}^{\mathrm{vertex}}) f\rangle =
\min_{\psi: p^{0}\psi = f}
\langle \psi, A_{n+1}^{\mathrm{vertex}} \psi\rangle =
\langle H_{n+1} f, A_{n+1}^{\mathrm{vertex}} H_{n+1} f\rangle$
for every $f \in X_n^{0}$. This is standard: for
$\psi = (f, f_I)$ in block form,
$\langle \psi, A_{n+1}^{\mathrm{vertex}} \psi\rangle =
\langle f, A_{BB} f\rangle + 2\langle f_I, A_{IB} f\rangle +
\langle f_I, A_{II} f_I\rangle$ is minimised in $f_I$ at
$f_I^{\star} = -A_{II}^{-1} A_{IB} f$ (since $A_{II}$ is positive
definite), with minimum value
$\langle f, (A_{BB} - A_{BI} A_{II}^{-1} A_{IB}) f\rangle$, and
$H_{n+1} f := (f, f_I^{\star})$ by definition of harmonic extension.

\smallskip
\textbf{Step 2.} Let $f \in X_n^{0}$ and let $\psi \in X_{n+1}^{0}$
satisfy $p^{0} \psi = f$. Apply the coboundary refinement relation
\eqref{eq:cobref}:
$p^{1} d_{n+1} \psi = \tfrac12 d_n p^{0} \psi = \tfrac12 d_n f$,
which is fixed independently of $\psi$ (subject to
$p^{0} \psi = f$). The averaging-by-$\tfrac12$ structure of $p^{1}$
on a 2-element fibre yields the calibration identity
\begin{equation}\label{eq:p1-calib}
   p^{1} (p^{1})^{*} = \tfrac12 \, I_{X_n^{1}},
\end{equation}
which is verified directly: for any $\omega \in X_n^{1}$ and any
level-$n$ edge $e$ with two children $e_1, e_2$, one has
$((p^{1})^{*}\omega)(e_i) = \tfrac12 \omega(e)$ (the unique adjoint
under standard $\ell^{2}$ inner products) and
$(p^{1}(p^{1})^{*}\omega)(e) = \tfrac12 \cdot \tfrac12 \omega(e) +
\tfrac12 \cdot \tfrac12 \omega(e) = \tfrac12 \omega(e)$.
Hence the orthogonal projection of any
$\xi \in X_{n+1}^{1}$ onto $\mathrm{range}((p^{1})^{*})$ is
$P \xi = (p^{1})^{*}\bigl((p^{1}(p^{1})^{*})^{-1} p^{1} \xi\bigr) =
2(p^{1})^{*} p^{1} \xi$, and $X_{n+1}^{1}$ decomposes orthogonally
as $\mathrm{range}((p^{1})^{*}) \oplus \ker p^{1}$. Apply this
decomposition to $\xi = d_{n+1} \psi$:
\begin{equation}\label{eq:dpsi-decomp}
   d_{n+1} \psi = (p^{1})^{*} d_n f \;+\; r_{\psi},
   \qquad r_{\psi} \in \ker p^{1},
\end{equation}
where the first term comes from
$P (d_{n+1} \psi) = 2 (p^{1})^{*} (p^{1} d_{n+1} \psi) =
2 (p^{1})^{*} (\tfrac12 d_n f) = (p^{1})^{*} d_n f$. Compute the
norm using \eqref{eq:p1-calib}:
\begin{align*}
\|(p^{1})^{*} d_n f\|^{2}
   = \langle (p^{1})^{*} d_n f, (p^{1})^{*} d_n f\rangle
   = \langle d_n f, p^{1} (p^{1})^{*} d_n f\rangle
   = \tfrac12 \|d_n f\|^{2}
   = \tfrac12 \langle f, A_n^{\mathrm{vertex}} f\rangle.
\end{align*}
By orthogonality of the decomposition \eqref{eq:dpsi-decomp},
\begin{equation}\label{eq:psi-energy-lb}
\langle \psi, A_{n+1}^{\mathrm{vertex}} \psi\rangle
   = \|d_{n+1} \psi\|^{2}
   = \tfrac12 \langle f, A_n^{\mathrm{vertex}} f\rangle + \|r_{\psi}\|^{2}
   \;\ge\; \tfrac12 \langle f, A_n^{\mathrm{vertex}} f\rangle.
\end{equation}

\smallskip
\textbf{Step 3.}
We exhibit one specific $\psi$ with $r_{\psi} = 0$, hence equality
in \eqref{eq:psi-energy-lb}: the harmonic extension itself. By the
edge-midpoint hypothesis of the theorem statement, every interior
vertex $m_e \in I$ has degree exactly~$2$ in $G_{n+1}^{\bullet}$,
its only neighbours being the two endpoints $u, v$ of its parent
edge $e \in E(G_n^{\bullet})$; hence the harmonic-extension formula
$f_I = -A_{II}^{-1} A_{IB} f$ at $m_e$ reduces to the elementary
two-neighbour average
$(H_{n+1} f)(m_e) = \tfrac12 (f(u) + f(v))$. Compute
$d_{n+1}(H_{n+1} f)$ at the two child edges of~$e$:
\begin{align*}
   (d_{n+1}(H_{n+1} f))(u, m_e)
   & = (H_{n+1} f)(m_e) - (H_{n+1} f)(u)
   = \tfrac12 (f(u) + f(v)) - f(u)
   = \tfrac12 (f(v) - f(u))
   = \tfrac12 (d_n f)(e), \\
   (d_{n+1}(H_{n+1} f))(m_e, v)
   & = (H_{n+1} f)(v) - (H_{n+1} f)(m_e)
   = f(v) - \tfrac12 (f(u) + f(v))
   = \tfrac12 (f(v) - f(u))
   = \tfrac12 (d_n f)(e).
\end{align*}
On the other hand,
$((p^{1})^{*} d_n f)(u, m_e) = \tfrac12 (d_n f)(e) =
((p^{1})^{*} d_n f)(m_e, v)$ by definition of $(p^{1})^{*}$ on a
two-element fibre. Hence
$d_{n+1}(H_{n+1} f) = (p^{1})^{*} d_n f$ on every level-$(n+1)$ edge,
i.e., $r_{H_{n+1} f} = 0$.

\smallskip
Combining Steps 1--3:
\begin{align*}
\langle f, \Schur(A_{n+1}^{\mathrm{vertex}}) f\rangle
   = \langle H_{n+1} f, A_{n+1}^{\mathrm{vertex}} H_{n+1} f\rangle
   = \tfrac12 \langle f, A_n^{\mathrm{vertex}} f\rangle
   + \|r_{H_{n+1} f}\|^{2}
   = \tfrac12 \langle f, A_n^{\mathrm{vertex}} f\rangle.
\end{align*}
This holds for every $f \in X_n^{0}$. Since
$\Schur(A_{n+1}^{\mathrm{vertex}})$ and $A_n^{\mathrm{vertex}}$ are
both symmetric, polarisation of the quadratic form gives the
operator identity \eqref{eq:schur-half}. The scaled form
\eqref{eq:schur-tilde} follows by linearity: $\Schur$ is a quadratic
operation but its dependence on the overall scale of the input
matrix is linear in that scale, so
$\Schur(\widetilde A_{n+1}) = \Schur(2^{n+1} A_{n+1}^{\mathrm{vertex}}) =
2^{n+1} \Schur(A_{n+1}^{\mathrm{vertex}}) =
2^{n+1} \cdot \tfrac12 A_n^{\mathrm{vertex}} = 2^n A_n^{\mathrm{vertex}} = \widetilde A_n$.
\end{proof}

\begin{remark}[Numerical verification]\label{rem:numerics}
Theorem~\ref{thm:schur} is verified to exact rational precision in
\S\ref{sec:worked} for two finite refinement instances: the single
edge with one-midpoint subdivision (Tower~I,
$|V(G_0^{\bullet})| = 2$, $|V(G_1^{\bullet})| = 3$) and the triangle
with edge-midpoint subdivision (Tower~II,
$|V(G_0^{\bullet})| = 3$, $|V(G_1^{\bullet})| = 6$). The
numerical witness is
\texttt{papers/cascade-refinement-compat/scripts/explore\_defect.py};
all linear algebra is exact \texttt{Fraction} arithmetic, with no
floating-point error.
\end{remark}

\section{Energy intertwining and discharge of (O3)}\label{sec:O3}

We now apply Theorem~\ref{thm:schur} to the cascade-mechanism
consumer-API obligations. The cascade-mechanism residual functional
in this paper's setting is
$\widetilde R_n(\phi) := \tfrac12 \langle \phi, \widetilde A_n \phi\rangle$.

\begin{theorem}[Energy intertwining; discharge of (O3)]\label{thm:O3-discharge}
In the abstract scalar pure-midpoint refinement model of
Definition~\ref{def:abstract-model}, for every
$\psi \in X_{n+1}^{0}$,
\begin{equation}\label{eq:O3}
   \widetilde R_n\bigl(p^{0}_{n+1, n} \psi\bigr)
   \;\le\;
   \widetilde R_{n+1}(\psi),
\end{equation}
with equality if and only if $\psi = H_{n+1}\bigl(p^{0}_{n+1, n} \psi\bigr)$
(i.e.\ $\psi$ is the harmonic extension of its boundary).
\end{theorem}

\begin{proof}
Decompose $\psi = H_{n+1}(p^{0}_{n+1, n} \psi) + \psi^{\perp}$, where
$\psi^{\perp} := \psi - H_{n+1}(p^{0}_{n+1, n} \psi)$ vanishes on
$B = V(G_n^{\bullet})$ (since
$p^{0}_{n+1, n} H_{n+1}(p^{0}_{n+1, n} \psi) = p^{0}_{n+1, n} \psi$).
The energy quadratic form decomposes as
\begin{align*}
\langle \psi, A_{n+1}^{\mathrm{vertex}} \psi\rangle
  & = \langle H_{n+1}(p^{0} \psi), A_{n+1}^{\mathrm{vertex}} H_{n+1}(p^{0} \psi)\rangle \\
  & \quad + 2 \langle H_{n+1}(p^{0} \psi), A_{n+1}^{\mathrm{vertex}} \psi^{\perp}\rangle
       + \langle \psi^{\perp}, A_{n+1}^{\mathrm{vertex}} \psi^{\perp}\rangle.
\end{align*}
The cross term vanishes: writing
$A_{n+1}^{\mathrm{vertex}}$ in block form \eqref{eq:block} and noting
that $\psi^{\perp}$ has zero $B$-component, one has
$A_{n+1}^{\mathrm{vertex}} \psi^{\perp} = (A_{BI} \psi^{\perp}_{I}, A_{II} \psi^{\perp}_{I})$,
and $H_{n+1}(p^{0} \psi) = (p^{0} \psi, -A_{II}^{-1} A_{IB} p^{0} \psi)$;
the cross-pairing computes to
$\langle p^{0} \psi, A_{BI} \psi^{\perp}_{I}\rangle
+ \langle -A_{II}^{-1} A_{IB} p^{0} \psi, A_{II} \psi^{\perp}_{I}\rangle
= \langle p^{0} \psi, A_{BI} \psi^{\perp}_{I}\rangle
- \langle p^{0} \psi, A_{BI} \psi^{\perp}_{I}\rangle = 0$.
The first surviving term equals
$\langle p^{0} \psi, \Schur(A_{n+1}^{\mathrm{vertex}}) p^{0} \psi\rangle =
\tfrac12 \langle p^{0} \psi, A_n^{\mathrm{vertex}} p^{0} \psi\rangle$
by Theorem~\ref{thm:schur}. The third surviving term is
$\langle \psi^{\perp}_{I}, A_{II} \psi^{\perp}_{I}\rangle \ge 0$
since $A_{II}$ is positive semi-definite (a principal submatrix of
the positive-semi-definite $A_{n+1}^{\mathrm{vertex}}$). Hence
\begin{equation*}
\langle \psi, A_{n+1}^{\mathrm{vertex}} \psi\rangle
\;\ge\;
\tfrac12 \langle p^{0} \psi, A_n^{\mathrm{vertex}} p^{0} \psi\rangle.
\end{equation*}
Multiplying both sides by $2^{n}$:
$2^{n} \langle \psi, A_{n+1}^{\mathrm{vertex}} \psi\rangle
\ge 2^{n-1} \langle p^{0} \psi, A_n^{\mathrm{vertex}} p^{0} \psi\rangle$,
which after dividing by $2$ and re-grouping factors becomes
$\langle \psi, \widetilde A_{n+1} \psi\rangle / 2 \ge
\langle p^{0} \psi, \widetilde A_n p^{0} \psi\rangle / 2$, i.e.\
$\widetilde R_{n+1}(\psi) \ge \widetilde R_n(p^{0} \psi)$, which is
\eqref{eq:O3}. Equality requires $\psi^{\perp} = 0$ (since $A_{II}$
is positive definite by invertibility), i.e.\ $\psi$ is the harmonic
extension of its boundary.
\end{proof}

\begin{corollary}[(O3) is closed in the abstract pure-midpoint model]\label{cor:O3-closed}
In the abstract scalar pure-midpoint refinement model of
Definition~\ref{def:abstract-model}, the cascade-mechanism
residual-monotonicity hypothesis \cite{CascadeMechanism}~(O3) holds
for the scaled curvature residual
$\widetilde R_n(\phi) = \tfrac12 \langle \phi, \widetilde A_n \phi\rangle$,
where $\widetilde A_n = 2^n A_n^{\mathrm{vertex}}$ and
$A_n^{\mathrm{vertex}} = d_n^{*} d_n$ is the vertex
graph Laplacian. Lifting this discharge to the full Schl\"{a}fli
P3/P4 tower of \cite{CascadeClosureDynamics} (which adds
face-centroid vertices and centroid--midpoint edges absent from
Definition~\ref{def:abstract-model}) is open and named in
\S\ref{sec:scope-and-gap}.
\end{corollary}

\section{Strict operator-level refinement-compatibility fails}\label{sec:strict-fails}

Theorem~\ref{thm:schur} closes the energy-level refinement question.
The strict operator-level identity
$p^{0}_{n+1, n} \widetilde A_{n+1} = \widetilde A_n p^{0}_{n+1, n}$
remains false, with explicit defect.

\begin{proposition}[Strict operator intertwining fails; kernel of the defect]\label{prop:strict-fails}
In the abstract scalar pure-midpoint refinement model of
Definition~\ref{def:abstract-model}, the strict identity
$p^{0}_{n+1, n} \widetilde A_{n+1} = \widetilde A_n p^{0}_{n+1, n}$ does
not hold in general, and equivalently
$p^{0}_{n+1, n} A_{n+1}^{\mathrm{vertex}} \ne \tfrac12 A_n^{\mathrm{vertex}} p^{0}_{n+1, n}$
as operators $X_{n+1}^{0} \to X_n^{0}$. The defect operator
\begin{equation}\label{eq:defect}
   D_n := p^{0}_{n+1, n} A_{n+1}^{\mathrm{vertex}} - \tfrac12 A_n^{\mathrm{vertex}} p^{0}_{n+1, n}
\end{equation}
satisfies the following:
\begin{enumerate}[label=\textup{(\roman*)}, leftmargin=2.4em, nosep, topsep=2pt]
\item $D_n \,\psi = 0$ whenever $\psi \in \mathrm{range}(H_{n+1})$,
  i.e.\ whenever $\psi$ is the harmonic extension of its own
  boundary $p^{0} \psi$;
\item Writing $\delta(\psi) \in \R^{E_n}$ for the per-edge midpoint
  deficit
  $\delta_e(\psi) := \tfrac12(\psi_u + \psi_v) - \psi_{m_e}$
  (with $e = \{u, v\} \in E_n$ and $m_e$ the midpoint inserted on
  $e$), one has $(D_n \psi)_v = \sum_{e \ni v} \delta_e(\psi)$ for
  every $v \in V_n$. Equivalently, $D_n \psi = B_n \delta(\psi)$
  where $B_n \in \{0, 1\}^{V_n \times E_n}$ is the unsigned
  vertex--edge incidence matrix of $G_n$. Consequently
  \begin{equation}\label{eq:kerDn}
     \ker D_n
     \;=\; \mathrm{range}(H_{n+1}) + \bigl\{\psi : \delta(\psi) \in \ker B_n,\ \psi_B = 0\bigr\},
  \end{equation}
  with the two summands intersecting trivially. The kernel
  $\ker B_n$ is non-trivial precisely when $G_n$ admits an
  alternating-sign edge-weighting along a closed walk of even
  length; in particular $\ker B_n = 0$ when $G_n$ is a forest, and
  $\ker D_n = \mathrm{range}(H_{n+1})$ in that case. The smallest
  counterexample is $G_n = C_4$ (4-cycle), where $\dim \ker B_n = 1$
  and $\ker D_n$ is strictly larger than
  $\mathrm{range}(H_{n+1})$;
\item $D_n$ does \emph{not} vanish on the level-$n$-supported
  subspace $\{\psi : \psi_I = 0\}$ in general (and not at all when
  $G_n$ has at least one edge).
\end{enumerate}
\end{proposition}

\begin{proof}
\textbf{Per-vertex formula.} Fix $v \in V_n$. By
Definition~\ref{def:abstract-model}~\textup{(M2)}, the level-$(n+1)$
neighbours of $v$ are exactly the midpoints
$\{m_e : e \in E_n,\, e \ni v\}$ of the level-$n$ edges incident to
$v$, and $v$ has the same level-$(n+1)$ degree
$\deg_{G_{n+1}}(v) = \deg_{G_n}(v)$. Hence
\[
   (A_{n+1}^{\mathrm{vertex}} \psi)(v)
   \;=\; \sum_{e \ni v} \bigl(\psi(v) - \psi(m_e)\bigr)
   \quad\text{and}\quad
   (A_n^{\mathrm{vertex}} \psi_B)(v)
   \;=\; \sum_{e = \{v, u\} \ni v} \bigl(\psi(v) - \psi(u)\bigr),
\]
so for every $\psi \in X_{n+1}^{0}$ with $\psi_B = p^{0} \psi$,
\begin{align*}
   (D_n \psi)(v)
   &= (A_{n+1}^{\mathrm{vertex}} \psi)(v)
      - \tfrac12 (A_n^{\mathrm{vertex}} \psi_B)(v) \\
   &= \sum_{e \ni v} \bigl(\psi(v) - \psi(m_e)\bigr)
      - \tfrac12 \sum_{e = \{v, u\}} \bigl(\psi(v) - \psi(u)\bigr) \\
   &= \sum_{e \ni v} \Bigl[\tfrac12(\psi(u) + \psi(v)) - \psi(m_e)\Bigr]
   \;=\; \sum_{e \ni v} \delta_e(\psi)
   \;=\; (B_n \delta(\psi))(v).
\end{align*}
This proves the formula in the proposition statement.

\smallskip
\textbf{(i)} If $\psi = H_{n+1} f$ is the harmonic extension of
$f \in X_n^{0}$, the harmonic-equation
$(A_{II})_{m_e, m_e} \psi(m_e) + (A_{IB})_{m_e, \cdot} f = 0$
collapses, since $m_e$ has level-$(n+1)$ neighbours
$\{u, v\}$ (so $A_{II}$ is diagonal with entries $2$ on midpoints,
$A_{IB}$ has rows $-e_u^{T} - e_v^{T}$ on midpoint $m_e$), to
$2 \psi(m_e) = f(u) + f(v)$, i.e.\
$\psi(m_e) = \tfrac12(f(u) + f(v))$. Hence
$\delta_e(\psi) = \tfrac12(f(u) + f(v)) - \psi(m_e) = 0$ for all
$e \in E_n$, and the per-vertex formula gives
$D_n \psi = B_n \cdot 0 = 0$.

\smallskip
\textbf{(ii)} From the per-vertex formula above, $D_n \psi = 0$ iff
$B_n \delta(\psi) = 0$ iff $\delta(\psi) \in \ker B_n$. The map
$\psi \mapsto \delta(\psi)$ is surjective
$X_{n+1}^{0} \twoheadrightarrow \R^{E_n}$ (any prescribed
$\delta \in \R^{E_n}$ is realised by $\psi_B := 0$ and
$\psi_{m_e} := -\delta_e$), and its kernel is exactly
$\mathrm{range}(H_{n+1})$ by the harmonic-extension equation
$\psi(m_e) = \tfrac12(\psi(u) + \psi(v))$ in part~(i). Hence the
short exact sequence
$0 \to \mathrm{range}(H_{n+1}) \to \ker D_n \to \ker B_n \to 0$
splits via the section $\delta \mapsto \psi$ defined by $\psi_B = 0$,
$\psi_{m_e} = -\delta_e$, giving \eqref{eq:kerDn}. The non-triviality
of $\ker B_n$ on closed even-length walks is the standard fact that
the unsigned incidence matrix has $\ker B_n^{T} B_n = \ker B_n$ with
$B_n^{T} B_n = 2 I_{E_n} + A^{\mathrm{adj}}_{\mathrm{line}}(G_n)$ on
the line graph, whose null space is generated by alternating-sign
$\pm 1$ patterns along even closed walks (and is trivial on
trees / forests). For $G_n = C_4$, an explicit kernel generator is
$\delta = (1, -1, 1, -1)$ on the four edges in cyclic order.

\smallskip
\textbf{(iii)} For any $\phi \in X_{n+1}^{0}$ with $\phi_I = 0$ and
edge $e = \{u, v\} \in E_n$, $\delta_e(\phi) = \tfrac12(\phi(u) + \phi(v))$.
Take any $\phi$ with $\phi(u) \ne -\phi(v)$ on at least one edge
(possible whenever $G_n$ has an edge): then $\delta(\phi) \ne 0$,
and the per-vertex formula gives $D_n \phi = B_n \delta(\phi)$ with
the right-hand side non-zero at any vertex incident to an edge $e$
where $\delta_e(\phi) \ne 0$ and not cancelled by sums at that
vertex. The explicit defect $D_n \phi$ is computed for two concrete
towers in \S\ref{sec:worked}; in both, $D_n |_{\{\phi_I = 0\}}$ is
non-zero.
\end{proof}

\begin{remark}[Geometric interpretation]\label{rem:harmonic-zero}
The conclusion (i) of Proposition~\ref{prop:strict-fails} is the
operator-level shadow of Theorem~\ref{thm:schur}: harmonic
extensions saturate the boundary Laplacian action under $p^{0}$ at
the strict scaled level. Conclusion (ii) records the genuinely
finer fact that, on parent graphs $G_n$ with non-trivial unsigned
incidence kernel (the smallest example is $G_n = C_4$), the kernel
of $D_n$ is strictly larger than $\mathrm{range}(H_{n+1})$: there
exist non-harmonic level-$(n+1)$ states whose midpoint deficits
form a non-zero edge-cocycle in $\ker B_n$ but whose total at every
boundary vertex still cancels. The defect $D_n$ measures the
boundary-vertex sum of midpoint deficits
$\delta_e(\psi) = \tfrac12(\psi_u + \psi_v) - \psi_{m_e}$ along
incident edges, not the per-edge midpoint deficit itself.
\end{remark}

\begin{remark}[Energy vs.\ flow level]
The Schur identity Theorem~\ref{thm:schur} and its energy
consequence Theorem~\ref{thm:O3-discharge} hold at the
\emph{quadratic-form / boundary-energy} level. Strict
\emph{operator}-level intertwining
$p^{0} \widetilde A_{n+1} = \widetilde A_n p^{0}$ fails by
Proposition~\ref{prop:strict-fails}, and consequently the strict
\emph{flow}-level intertwining
$p^{0} \circ e^{-t \widetilde A_{n+1}} = e^{-t \widetilde A_n} \circ p^{0}$
also fails. Cascade-mechanism's obligation~(O2), as stated, is at
the flow level. The natural source-supported reformulation of
(O2) is the Schur-energy version, which Theorem~\ref{thm:schur}
discharges. In the cascade-mechanism reformulation table of
\S\ref{sec:cascade-mechanism-update} we record this as a
\emph{partial} discharge of (O2).
\end{remark}

\section{Kernel of $A_n^{\mathrm{vertex}}$ and (O1)}\label{sec:O1}

For completeness we record the status of cascade-mechanism's
obligation~(O1) (top-rung residual-zero fixed point and convergence)
in the abstract scalar pure-midpoint refinement model of
Definition~\ref{def:abstract-model}, applied at the top level
$n = N$.

\begin{proposition}[Kernel and convergence of $\widetilde A_N$]\label{prop:O1-status}
For any finite level $N$, $\widetilde A_N$ is positive semi-definite
on $X_N^{0}$ and its kernel
$\ker \widetilde A_N = \ker A_N^{\mathrm{vertex}}$ consists exactly of
the locally-constant functions on the connected components of
$G_N^{\bullet}$. If $G_N^{\bullet}$ is connected,
$\dim \ker \widetilde A_N = 1$ (constants only).

For every initial datum $\phi^{(0)} \in X_N^{0}$, the gradient flow
$\Psi^{t}_N := e^{-t \widetilde A_N}$ converges as $t \to \infty$ to
the orthogonal projection of $\phi^{(0)}$ onto
$\ker \widetilde A_N$. This is automatic in finite dimensions by
spectral decomposition.

In particular, the trivial state $\phi^{\star}_N = 0$ is always a
residual-zero fixed point, and the cascade-mechanism obligation~(O1)
``nontrivial-selected-state existence'' reduces to whether
$\ker \widetilde A_N \ne \{0\}$ (yes, in the connected case the
kernel is one-dimensional) and whether the projection of the
initial datum onto the kernel is non-zero (which is generic but
not automatic).
\end{proposition}

\begin{proof}
$A_N^{\mathrm{vertex}} = d_N^{*} d_N$ is positive semi-definite by
construction, with kernel $\ker d_N$, the locally-constant
functions on the connected components of $G_N^{\bullet}$. The
convergence statement for $e^{-t A_N^{\mathrm{vertex}}}$ is the
standard finite-dimensional functional-calculus result for a
self-adjoint positive-semi-definite operator. Multiplying by the
positive scalar $2^N$ does not change the kernel or the convergence
behaviour.
\end{proof}

\begin{remark}[Selection rule is not part of cascade closure]
\cite{CascadeMechanism}~\S5 explicitly states that selection of
the top-rung fixed point is not provided by cascade closure and is
the successor paper's own task. The kernel description above
records what cascade closure inherits at the top rung; the
selection rule by which a particular non-zero element of
$\ker \widetilde A_N$ is preferred over the zero state is a
question for downstream papers, and outside the scope of the
present source-side discharge.
\end{remark}

\section{Worked finite refinement towers}\label{sec:worked}

This section gives explicit matrix-level computations on two
concrete finite refinement towers, verifying
Theorems~\ref{thm:schur} and~\ref{thm:O3-discharge} and
Proposition~\ref{prop:strict-fails} to exact rational precision.
The numerical witness is the script
\texttt{papers/cascade-refinement-compat/scripts/explore\_defect.py};
its output is preserved in
\texttt{papers/cascade-refinement-compat/repro/expected\_outputs/}.

\subsection{Tower~I: single edge $\to$ one-midpoint subdivision}

Let $G_0^{\bullet}$ be the graph on
$V(G_0^{\bullet}) = \{u, v\}$ with the single edge $e = (u, v)$.
Refine by inserting the midpoint $m$ of $e$, splitting $e$ into
two new edges $(u, m)$ and $(m, v)$:
$V(G_1^{\bullet}) = \{u, m, v\}$,
$E(G_1^{\bullet}) = \{(u, m), (m, v)\}$. The discrete coboundaries
in oriented form are
$d_0 = \begin{pmatrix} -1 & 1 \end{pmatrix}$ on $X_0^{0} = \R^{2}$,
and
$d_1 = \begin{pmatrix} -1 & 1 & 0 \\ 0 & -1 & 1 \end{pmatrix}$ on
$X_1^{0} = \R^{3}$. The bonding maps are
$p^{0}_{1, 0} = \begin{pmatrix} 1 & 0 & 0 \\ 0 & 0 & 1 \end{pmatrix}$
and
$p^{1}_{1, 0} = \begin{pmatrix} \tfrac12 & \tfrac12 \end{pmatrix}$.
The vertex Laplacians are
\[
   A_0^{\mathrm{vertex}} =
   \begin{pmatrix} 1 & -1 \\ -1 & 1 \end{pmatrix},
   \qquad
   A_1^{\mathrm{vertex}} =
   \begin{pmatrix} 1 & -1 & 0 \\ -1 & 2 & -1 \\ 0 & -1 & 1 \end{pmatrix}.
\]
With the partition $B = \{u, v\}$, $I = \{m\}$, the block
decomposition yields
$A_{BB} = \begin{pmatrix} 1 & 0 \\ 0 & 1 \end{pmatrix}$,
$A_{BI} = \begin{pmatrix} -1 \\ -1 \end{pmatrix}$,
$A_{IB} = \begin{pmatrix} -1 & -1 \end{pmatrix}$,
$A_{II} = \begin{pmatrix} 2 \end{pmatrix}$. The Schur complement is
\[
\Schur(A_1^{\mathrm{vertex}}) =
   A_{BB} - A_{BI} A_{II}^{-1} A_{IB}
   = \begin{pmatrix} 1 & 0 \\ 0 & 1 \end{pmatrix}
   - \tfrac12 \begin{pmatrix} 1 & 1 \\ 1 & 1 \end{pmatrix}
   = \begin{pmatrix} \tfrac12 & -\tfrac12 \\ -\tfrac12 & \tfrac12 \end{pmatrix}
   = \tfrac12 A_0^{\mathrm{vertex}}.
\]
This is exactly Theorem~\ref{thm:schur}.

The defect of strict refinement-compatibility is
\[
   D_0
   := p^{0}_{1, 0} A_1^{\mathrm{vertex}} - \tfrac12 A_0^{\mathrm{vertex}} p^{0}_{1, 0}
   = \begin{pmatrix} 1 & -1 & 0 \\ 0 & -1 & 1 \end{pmatrix}
   - \begin{pmatrix} \tfrac12 & 0 & -\tfrac12 \\ -\tfrac12 & 0 & \tfrac12 \end{pmatrix}
   = \begin{pmatrix} \tfrac12 & -1 & \tfrac12 \\ \tfrac12 & -1 & \tfrac12 \end{pmatrix}.
\]
Both rows of $D_0$ act on $\phi = (\phi(u), \phi(m), \phi(v))$ as
$\tfrac12 \phi(u) - \phi(m) + \tfrac12 \phi(v) =
\tfrac12 (\phi(u) + \phi(v)) - \phi(m)$, the ``midpoint deficit''
of $\phi$ at $m$. Tower~I has a single edge with no cycles, so
the unsigned vertex--edge incidence matrix
$B_0 \in \{0, 1\}^{2 \times 1}$ has trivial kernel, and the
incidence characterisation in
Proposition~\ref{prop:strict-fails}(ii) collapses to: $D_0 \phi = 0$
iff the midpoint deficit at $m$ is zero, iff
$\phi(m) = \tfrac12 (\phi(u) + \phi(v))$, iff $\phi$ is the harmonic
extension of its boundary. So here the kernel of $D_0$
coincides with $\mathrm{range}(H_1)$. Furthermore, $D_0$ does
\emph{not} vanish on $\{\phi : \phi(m) = 0\}$ (which gives
$D_0 \phi = (\tfrac12 (\phi(u) + \phi(v)),
\tfrac12 (\phi(u) + \phi(v)))$, generically non-zero), confirming
Proposition~\ref{prop:strict-fails}(iii).

\subsection{Tower~II: triangle $\to$ edge-midpoint subdivision}

Let $G_0^{\bullet}$ be the 3-cycle on
$V(G_0^{\bullet}) = \{0, 1, 2\}$ with edges
$e_{01} = (0, 1)$, $e_{12} = (1, 2)$, $e_{20} = (2, 0)$. Refine by
inserting one midpoint per edge, $m_{01}, m_{12}, m_{20}$, giving
$V(G_1^{\bullet}) = \{0, 1, 2, m_{01}, m_{12}, m_{20}\}$
($|V(G_1^{\bullet})| = 6$) and the six new edges
$(0, m_{01}), (m_{01}, 1), (1, m_{12}), (m_{12}, 2), (2, m_{20}), (m_{20}, 0)$.
With $B = \{0, 1, 2\}$, $I = \{m_{01}, m_{12}, m_{20}\}$ and the
identification of bonding maps
$p^{0}_{1, 0}$, $p^{1}_{1, 0}$ analogously to Tower~I, direct
computation (verified by
\texttt{explore\_defect.py}) yields
\[
   A_0^{\mathrm{vertex}} =
   \begin{pmatrix} 2 & -1 & -1 \\ -1 & 2 & -1 \\ -1 & -1 & 2 \end{pmatrix}
\]
and the Schur complement
\[
   \Schur(A_1^{\mathrm{vertex}}) =
   \begin{pmatrix} 1 & -\tfrac12 & -\tfrac12 \\ -\tfrac12 & 1 & -\tfrac12 \\ -\tfrac12 & -\tfrac12 & 1 \end{pmatrix}
   = \tfrac12 A_0^{\mathrm{vertex}},
\]
again confirming Theorem~\ref{thm:schur}. Eigenvalues of the
defect operator
$\Schur(A_1^{\mathrm{vertex}}) - A_0^{\mathrm{vertex}}$ on this
tower are $\{0, -\tfrac32, -\tfrac32\}$, confirming the
ordering $\Schur \le A_0^{\mathrm{vertex}}$ (i.e.\ the
unscaled Schur complement under-estimates the level-$0$
Laplacian; rescaling by $2^n$ at level $n$ exactly compensates).

\section{Cascade-mechanism update}\label{sec:cascade-mechanism-update}

The present paper updates the worked-invocation status table of
\cite{CascadeMechanism}~\S3.5 in the abstract scalar pure-midpoint
refinement model of Definition~\ref{def:abstract-model}, which
\S\ref{sec:scope-and-gap} relates to the imported P3/P4
finite-level refinement tower of \cite{CascadeClosureDynamics}: any
update lifts to that tower under the named hypotheses
\textup{(L1)}--\textup{(L3)} of \S\ref{sec:scope-and-gap} and
otherwise stays inside the abstract model. The consumer-API
obligations themselves (\cite{CascadeMechanism}~\S3.4) are
unchanged; the cascade-closure mechanism, Definition, Proposition,
and Corollary of \cite{CascadeMechanism} are unchanged.

\begin{longtable}{>{\raggedright\arraybackslash}p{0.20\textwidth}>{\raggedright\arraybackslash}p{0.40\textwidth}>{\raggedright\arraybackslash}p{0.32\textwidth}}
\caption{Updated status of cascade-mechanism consumer-API obligations
in the abstract scalar pure-midpoint refinement model
(Definition~\ref{def:abstract-model}) with the scaled curvature
operator $\widetilde A_n = 2^n A_n^{\mathrm{vertex}}$. Lifting any
row to the full P3/P4 tower of \cite{CascadeClosureDynamics}
requires the named hypotheses
\textup{(L1)}--\textup{(L3)} of \S\ref{sec:scope-and-gap}.}\label{tab:status-update}\\
\toprule
\textbf{Obligation} & \textbf{Status (this paper)} & \textbf{Where established} \\
\midrule
\endfirsthead
\toprule
\textbf{Obligation} & \textbf{Status (this paper)} & \textbf{Where established} \\
\midrule
\endhead
(O0) Standing analytic / admissibility hypotheses
& \textbf{Closed.} Finite-dimensional Hilbert structure of $X_n^{0}$, $C^1$ bonding maps, well-posed flow $e^{-t \widetilde A_n}$.
& \cite{CascadeRefinementSpaces}; recap in \S\ref{sec:intro}. \\
(O1) Top-rung fixed point and convergence
& \textbf{Closed at the trivial level} (zero state always residual-zero); \textbf{nontrivial selection} reduces to: kernel of $\widetilde A_N$ is the constants on connected components of $G_N^{\bullet}$, and convergence is automatic by spectral decomposition. The selection rule itself is out of scope per \cite{CascadeMechanism}~\S5.
& Proposition~\ref{prop:O1-status}. \\
(O2) Flow intertwining
& \textbf{Closed at the Schur / quadratic-form level only.} Strict operator-level identity $p^{0} \widetilde A_{n+1} = \widetilde A_n p^{0}$ remains false (Proposition~\ref{prop:strict-fails}); strict flow-level identity also remains false. The boundary-energy form satisfies $\Schur(\widetilde A_{n+1}) = \widetilde A_n$ exactly (Theorem~\ref{thm:schur}).
& Theorem~\ref{thm:schur}, Proposition~\ref{prop:strict-fails}. \\
(O3) Residual monotonicity
& \textbf{Closed.} $\widetilde R_n(p^{0} \psi) \le \widetilde R_{n+1}(\psi)$ for every $\psi \in X_{n+1}^{0}$, with equality iff $\psi$ is the harmonic extension of its boundary.
& Theorem~\ref{thm:O3-discharge}, Corollary~\ref{cor:O3-closed}. \\
\bottomrule
\end{longtable}

The discharge of (O3) and the partial discharge of (O2) are
non-trivial relative to \cite{CascadeMechanism}'s original status
report; the structural discovery is that the source's coboundary
factor-$\tfrac12$ relation \eqref{eq:cobref} lifts to the curvature
level via the Schur complement, with the scaled operator
$\widetilde A_n := 2^n A_n^{\mathrm{vertex}}$ as the natural
refinement-compatible object at the boundary-energy level.

\section{Discussion}\label{sec:discussion}

\paragraph*{What this paper does.}
It proves that, in the abstract scalar pure-midpoint refinement
model of Definition~\ref{def:abstract-model} and with the scaled
curvature operator $\widetilde A_n = 2^n A_n^{\mathrm{vertex}}$,
the cascade-mechanism residual-monotonicity obligation~(O3) holds
(Theorem~\ref{thm:O3-discharge}, Corollary~\ref{cor:O3-closed}) and
the flow-intertwining obligation~(O2) holds at the
boundary-energy / Schur level (Theorem~\ref{thm:schur}). The
mechanism by which both happen is the same: the coboundary
factor-$\tfrac12$ relation
$p^{1} d_{n+1} = \tfrac12 d_n p^{0}$, which holds verbatim in the
abstract model (and as a theorem in the source's full Schl\"{a}fli
setting), pulled forward through the boundary--interior
decomposition.

\paragraph*{What this paper does not do.}
It does not establish strict operator-level
$p^{0} \widetilde A_{n+1} = \widetilde A_n p^{0}$
intertwining; that identity is genuinely false in the abstract
model. The defect operator $D_n$ defined in
Proposition~\ref{prop:strict-fails} satisfies
$\mathrm{range}(H_{n+1}) \subseteq \ker D_n$, with equality when
the parent graph $G_n$ has no bipartite component, and $D_n$ is
non-zero on the level-$n$-supported sector
$\{\psi : \psi_I = 0\}$ in general. In particular, the strict
cascade-mechanism flow-intertwining hypothesis~(O2) at the
operator/flow level remains a strictly stronger condition than
what the abstract model supports; cascade-mechanism users invoking
the boundary-energy version of (O2) will find it discharged here,
but those requiring strict operator-level intertwining will find
the defect computed but not zero. The paper also does not lift any
of these results to the full Schl\"{a}fli P3/P4 tower of
\cite{CascadeClosureDynamics}: that lift is named in
\S\ref{sec:scope-and-gap} as hypotheses
\textup{(L1)}--\textup{(L3)} and is not attempted here. The paper
does not derive a selection rule for nontrivial elements of
$\ker \widetilde A_N$, in keeping with
\cite{CascadeMechanism}~\S5's explicit boundary.

\paragraph*{Implication for cascade-mechanism v2.1.}
The minimum cascade-mechanism update is a single sentence in
\cite{CascadeMechanism}~\S3.5 replacing the ``(O3) open'' line
with a citation of Corollary~\ref{cor:O3-closed} above, scoped to
the abstract pure-midpoint model of
Definition~\ref{def:abstract-model}, and a short remark recording
the partial discharge of (O2) at the Schur level under the same
scope. The cascade-mechanism consumer API itself
(\cite{CascadeMechanism}~\S3.4 (O0)--(O3)) does not change. The
reproduction bundle of \cite{CascadeMechanism} is unaffected. Any
update describing the full P3/P4 tower must be flagged conditional
on \textup{(L1)}--\textup{(L3)} of \S\ref{sec:scope-and-gap}.

\paragraph*{Open follow-ups.}
Proposition~\ref{prop:strict-fails}(ii) characterises $\ker D_n$ via
the unsigned vertex--edge incidence operator $B_n$. The natural
follow-up is the structure of $\ker D_n \setminus \mathrm{range}(H_{n+1})$
when $G_n$ has non-trivial $\ker B_n$ (e.g.\ $G_n = C_4$), and whether such
``incidence-cocycle'' kernel directions form an invariant sector
under the cascade-mechanism flow $e^{-t \widetilde A_{n+1}}$
(numerically they do not in general, since $\widetilde A_{n+1}$ does
not preserve the harmonic subspace). A separate question is
whether a non-linear cascade-flow modification (gradient flow on a
modified residual that absorbs the defect into the dynamics)
preserves the mechanism's structure while admitting strict
intertwining; this is out of scope for the present paper. The
P3/P4-tower lift under \textup{(L1)}--\textup{(L3)} is the most
load-bearing follow-up.

\section{Reproducibility}\label{sec:repro}

All computations in \S\S\ref{sec:schur}--\ref{sec:worked} are
verified to exact rational precision (\texttt{Fraction} arithmetic)
by the script
\texttt{papers/cascade-refinement-compat/scripts/explore\_defect.py}.
A reproducer with frozen reference outputs lives at
\texttt{papers/cascade-refinement-compat/repro/}; from the
repository root, one invocation
\begin{verbatim}
cd papers/cascade-refinement-compat/repro
bash run_audit.sh
\end{verbatim}
re-runs the verification on Towers~I and~II and diffs against the
frozen expected outputs. A clean run prints
\texttt{REFINEMENT COMPAT AUDIT PASSED}.

\bibliographystyle{plain}
\bibliography{references}

\end{document}
